% GENERATED FILE. Edit canonical lessons, then run npm run generate:books.
% canonical-source-hash: fnv1a64:4667bac32c6b3aab
% canonical-lessons: LA-C52-ubi, LA-C52-cur, LA-C52-quia, LA-C52-meus, LA-C52-tuus

\chapter{Ubi, Cur, Quia, Meus, Tuus}
\label{ch:la-interrogatio}
\hlchaptermodality{52}

\begin{chapteropening}
\textbf{By the end of this chapter:} \emph{I can ask where and why, answer with a because-clause, and say that a thing is mine or yours.}

The fifth word of the group, taken apart against the four before it.
\end{chapteropening}

\section[ubi]{ubi — the question the \emph{quōmodo} lesson left half-finished}

\label{lesson:LA-C52-ubi}



\begin{quote}
Six recalls before the new word, at growing distances. They are not decoration: each one is the retrieval window this book measures, and each names a word far enough back that saying it is work.

\begin{itemize}
\raggedright
  \item \emph{Recall:} say \emph{epistula} — \textbf{R1}, one lesson back, before it decays
  \item \emph{Recall:} say \emph{amīcus} — \textbf{R2}, the first real retrieval, five to fifteen lessons back
  \item \emph{Recall:} say \emph{nesciō} — \textbf{R3}, durable, twenty to sixty lessons back
  \item \emph{Recall:} say \emph{nesciō} — \textbf{R3}, durable again, a second word from the same distance
  \item \emph{Recall:} say \emph{tē volup est convēnisse} — \textbf{R4}, recognition at distance, eighty lessons back or more
  \item \emph{Recall:} say \emph{quid tibi nōmen est? — volup est convēnisse} — \textbf{R4}, recognition at distance, a second word from that far back
\end{itemize}
\end{quote}

\begin{cousinweb}
\textbf{ubi} — \textquotedblleft{}\textbf{where}.\textquotedblright{}

The \emph{quōmodo} lesson gave you \textquotedblleft{}how,\textquotedblright{} and stopped. This is the other half of the pair every beginner needs, and the corpus went fifty-one chapters without it.

\begin{quote}
\textbf{Ubi est liber?} — Where is the book? \textbf{Ibi.} — There. (Last chapter's word, finally with a question to answer.)
\end{quote}

Like \emph{quī} and \emph{quandō}, it is built on the \emph{\textbf{*kwo-}} stem. Latin also has \textbf{quō} (\textquotedblleft{}to where\textquotedblright{}) and \textbf{unde} (\textquotedblleft{}from where\textquotedblright{}) — Latin distinguishes three directions where English says only \emph{where}, and this is the resting one.

\textbf{Ubique} is \textquotedblleft{}everywhere,\textquotedblright{} and English \textbf{ubiquitous} is that word: being in every where at once. Spanish \emph{donde} and French \emph{où} both descend from \emph{ubi} with \emph{dē} in front of it.
\end{cousinweb}

\subsection*{Your turn}
Say these aloud:

\begin{itemize}
\raggedright
  \item \emph{ubi}
  \item every word this chapter has given you so far, in order
\end{itemize}

\begin{tcolorbox}[breakable,colback=teal!4,colframe=teal!35!black,title={Before you move on}]
Which earlier question word is \emph{ubi} the pair of? (\emph{\textbf{Quōmodo}}, how.) What does \emph{ubique} mean, and what English word is it? (\textbf{Everywhere}; \emph{\textbf{ubiquitous}}.)
\end{tcolorbox}
\section[cūr]{cūr — and the two words that answer it}

\label{lesson:LA-C52-cur}



\begin{quote}
Six recalls before the new word, at growing distances. They are not decoration: each one is the retrieval window this book measures, and each names a word far enough back that saying it is work.

\begin{itemize}
\raggedright
  \item \emph{Recall:} say \emph{ubi} — \textbf{R1}, one lesson back, before it decays
  \item \emph{Recall:} say \emph{pecūnia} — \textbf{R2}, the first real retrieval, five to fifteen lessons back
  \item \emph{Recall:} say \emph{male} — \textbf{R3}, durable, twenty to sixty lessons back
  \item \emph{Recall:} say \emph{capiō} — \textbf{R3}, durable again, a second word from the same distance
  \item \emph{Recall:} say \emph{tē volup est convēnisse} — \textbf{R4}, recognition at distance, eighty lessons back or more
  \item \emph{Recall:} say \emph{tū / vōs} — \textbf{R4}, recognition at distance, a second word from that far back
\end{itemize}
\end{quote}

\begin{cousinweb}
\textbf{cūr} — \textquotedblleft{}\textbf{why}.\textquotedblright{}

The third of this chapter's question words, and the one that opens the only question a learner cannot fake an answer to.

\begin{quote}
\textbf{Cūr nōn venīs?} — Why are you not coming? (with the \emph{nescio} lesson's rule.)
\end{quote}

\emph{Cūr} is an older \emph{quor}, the same \emph{\textbf{*kwo-}} stem again — \emph{quī}, \emph{quandō}, \emph{ubi}, \emph{cūr}, \emph{quid}, \emph{quis}, all one family. Recognise the \emph{qu-} and you can spot a Latin question word without knowing which one it is.

And like \emph{sed}, it \textbf{left no descendant}. Spanish asks \emph{¿por qué?}, Italian \emph{perché}, French \emph{pourquoi}, all rebuilt from \emph{per} or \emph{pro} plus \emph{quid} — \textquotedblleft{}for what.\textquotedblright{} Latin's own word for why did not survive its own daughters.
\end{cousinweb}

\subsection*{Your turn}
Say these aloud:

\begin{itemize}
\raggedright
  \item \emph{cūr}
  \item every word this chapter has given you so far, in order
\end{itemize}

\begin{tcolorbox}[breakable,colback=teal!4,colframe=teal!35!black,title={Before you move on}]
What do \emph{quī}, \emph{quandō}, \emph{ubi} and \emph{cūr} share? (\textbf{The \emph{*kwo-} stem} — Latin's question family.) Did any Romance language keep \emph{cūr}? (\textbf{No} — all rebuilt it from \emph{per} + \emph{quid}.)
\end{tcolorbox}
\section[quia]{quia — the answer cūr asks for}

\label{lesson:LA-C52-quia}



\begin{quote}
Six recalls before the new word, at growing distances. They are not decoration: each one is the retrieval window this book measures, and each names a word far enough back that saying it is work.

\begin{itemize}
\raggedright
  \item \emph{Recall:} say \emph{cūr} — \textbf{R1}, one lesson back, before it decays
  \item \emph{Recall:} say \emph{medicus} — \textbf{R2}, the first real retrieval, five to fifteen lessons back
  \item \emph{Recall:} say \emph{malus} — \textbf{R3}, durable, twenty to sixty lessons back
  \item \emph{Recall:} say \emph{rogō} — \textbf{R3}, durable again, a second word from the same distance
  \item \emph{Recall:} say \emph{tū / vōs} — \textbf{R4}, recognition at distance, eighty lessons back or more
  \item \emph{Recall:} say \emph{nihil est} — \textbf{R4}, recognition at distance, a second word from that far back
\end{itemize}
\end{quote}

\begin{cousinweb}
\textbf{quia} — \textquotedblleft{}\textbf{because}.\textquotedblright{}

The word that completes the pair. \emph{Cūr} asks; \emph{quia} answers, and it hangs a whole second statement onto the first:

\begin{quote}
\textbf{Cūr nōn venīs? Quia nōn possum.} — Why are you not coming? Because I cannot.
\end{quote}

That sentence uses four things this tranche has given you — \emph{cūr}, \emph{quia}, the \textbf{nōn} in front of a verb, and \textbf{possum} — and nothing else. It is the first question-and-reason exchange this book can build.

Latin has three words here and they are not quite the same: \textbf{quia} states a plain cause, \textbf{quod} states one the hearer already accepts, and \textbf{quoniam} (\textquotedblleft{}since now\textquotedblright{}) states one both parties take for granted. A reader meets all three. Start with \emph{quia}.
\end{cousinweb}

\subsection*{Your turn}
Say these aloud:

\begin{itemize}
\raggedright
  \item \emph{quia}
  \item every word this chapter has given you so far, in order
\end{itemize}

\begin{tcolorbox}[breakable,colback=teal!4,colframe=teal!35!black,title={Before you move on}]
Answer \emph{cūr nōn venīs?} using this chapter and the last two. (\emph{\textbf{Quia nōn possum}}.) How many Latin words mean \emph{because}? (\textbf{Three} — \emph{quia}, \emph{quod}, \emph{quoniam}.)
\end{tcolorbox}
\section[meus, mea, meum]{meus — the possessive the name lesson did without}

\label{lesson:LA-C52-meus}



\begin{quote}
Six recalls before the new word, at growing distances. They are not decoration: each one is the retrieval window this book measures, and each names a word far enough back that saying it is work.

\begin{itemize}
\raggedright
  \item \emph{Recall:} say \emph{quia} — \textbf{R1}, one lesson back, before it decays
  \item \emph{Recall:} say \emph{mīles} — \textbf{R2}, the first real retrieval, five to fifteen lessons back
  \item \emph{Recall:} say \emph{sed} — \textbf{R3}, durable, twenty to sixty lessons back
  \item \emph{Recall:} say \emph{adiuvō} — \textbf{R3}, durable again, a second word from the same distance
  \item \emph{Recall:} say \emph{tū / vōs} — \textbf{R4}, recognition at distance, eighty lessons back or more
  \item \emph{Recall:} say \emph{nihil est} — \textbf{R4}, recognition at distance, a second word from that far back
\end{itemize}
\end{quote}

\begin{cousinweb}
\textbf{meus, mea, meum} — \textquotedblleft{}\textbf{my}.\textquotedblright{}

The name-exchange lesson taught you to say \textbf{mihi nōmen est}, \textquotedblleft{}the name is \textbf{to me},\textquotedblright{} and called it the genuinely classical way to give your name. That is Latin's \emph{dative of possession}, and it is why this word waited so long: for a name, a Roman preferred the case to the adjective.

For everything else, the adjective is what you want, and it agrees with what is owned rather than with the owner:

\begin{quote}
\textbf{liber meus} — my book (masculine, so \emph{meus}) \textbf{epistula mea} — my letter (feminine, so \emph{mea})
\end{quote}

That is the opposite of English, where \emph{my} never changes. It matches French \emph{mon livre / ma lettre} and Spanish \emph{mi} only in part — Spanish flattened the gender out of \emph{mi} and kept it in \emph{mío}.
\end{cousinweb}

\subsection*{Your turn}
Say these aloud:

\begin{itemize}
\raggedright
  \item \emph{meus, mea, meum}
  \item every word this chapter has given you so far, in order
\end{itemize}

\begin{tcolorbox}[breakable,colback=teal!4,colframe=teal!35!black,title={Before you move on}]
What does \emph{meus} agree with, the owner or the thing owned? (\textbf{The thing owned}.) What did the name exchange use instead for a name? (\emph{\textbf{Mihi nōmen est}} — the dative of possession.)
\end{tcolorbox}
\section[tuus, tua, tuum]{tuus — and the pair of forms the tū / vōs lesson promised}

\label{lesson:LA-C52-tuus}



\begin{quote}
Six recalls before the new word, at growing distances. They are not decoration: each one is the retrieval window this book measures, and each names a word far enough back that saying it is work.

\begin{itemize}
\raggedright
  \item \emph{Recall:} say \emph{meus} — \textbf{R1}, one lesson back, before it decays
  \item \emph{Recall:} say \emph{epistula} — \textbf{R2}, the first real retrieval, five to fifteen lessons back
  \item \emph{Recall:} say \emph{aut} — \textbf{R3}, durable, twenty to sixty lessons back
  \item \emph{Recall:} say \emph{amō} — \textbf{R3}, durable again, a second word from the same distance
  \item \emph{Recall:} say \emph{nihil est} — \textbf{R4}, recognition at distance, eighty lessons back or more
  \item \emph{Recall:} say \emph{propediem tē vidēbō} — \textbf{R4}, recognition at distance, a second word from that far back
\end{itemize}
\end{quote}

\begin{cousinweb}
\textbf{tuus, tua, tuum} — \textquotedblleft{}\textbf{your},\textquotedblright{} to one person.

The \emph{tū / vōs} lesson made a point of it: in Classical Latin the difference is \textbf{number only}, with no politeness in it at all. This is \emph{tū}'s possessive, and it inherits that: \textbf{tuus} is not familiar or rude, it simply means the thing belongs to one person. For several, Latin says \textbf{vester}.

\begin{quote}
\textbf{Quid tibi nōmen est?} — the name question, with \emph{tibi}, the dative. \textbf{Liber tuus.} — Your book, with the adjective.
\end{quote}

\emph{Tibi} and \emph{tuus} are the same pronoun by two routes, which is worth hearing once: the case form and the adjective come from one word, and Latin lets you choose.

Spanish \emph{tu}, Italian \emph{tuo}, French \emph{ton} and Portuguese \emph{teu} all descend from this — and every one of them acquired the politeness dimension that Latin's had not got.
\end{cousinweb}

\subsection*{Your turn}
Say these aloud:

\begin{itemize}
\raggedright
  \item \emph{tuus, tua, tuum}
  \item every word this chapter has given you so far, in order
\end{itemize}

\begin{tcolorbox}[breakable,colback=teal!4,colframe=teal!35!black,title={Before you move on}]
Is \emph{tuus} familiar or polite? (\textbf{Neither} — Classical Latin marks number only.) Which earlier word is its case-form twin? (\emph{\textbf{Tibi}}.)
\end{tcolorbox}
